% !TeX program = lualatex
% Agregátor reálných otázek SZZ pro okruh M1. Jednotlivé otázky jsou v souborech q_*.tex.
% Sestavení: ./build.sh 01_spolecna_matematika/M1_diferencialni_integralni_pocet/priklady   (nebo cd .../priklady && latexmk)
\documentclass[11pt,a4paper]{article}
\usepackage{szzpriklady}
% Build běží z adresáře priklady/ (CWD), obrázky jsou v src/fig/.
\graphicspath{{src/fig/}{fig/}}

\title{Příklady otázek -- Základy diferenciálního a integrálního počtu\\[0.3em]\large Reálné otázky ze SZZ 2022--2026 (M1)}
\author{Jakub Klesa}
\date{Pokryté předměty: NMAI054 (Matematická analýza 1)}

\begin{document}
\maketitle
\tableofcontents
\bigskip

\begin{poznamka}
Tento dokument obsahuje reálné otázky z minulých termínů SZZ (2022--2026) zařazené do
okruhu M1. U otázek s obrázkem je grafika oříznutá z originálního PDF (viz odkaz na
zdroj). Text, číselná data a tabulky jsou přepsané věrně.
Zadání i oficiální náčrty řešení pocházejí ze sad zveřejněných na webu MFF UK:\par
{\raggedright\url{https://www.mff.cuni.cz/cs/studenti/bakalarske-studium/statni-zaverecne-zkousky/bakalarske-statni-zkousky-studijniho-programu-informatika}\par}
Náčrty označené jako doplněné jsou vlastní, oficiální sada je k~dané otázce neobsahuje.
\end{poznamka}
\bigskip

% === Otázky (chronologicky). Doplň \input řádky a soubory q_*.tex do src/. ===
% Příklad: \input{q_2024-leto-28_bayesovske-site}
% ZACATEK-OTAZEK
\input{src/q_2022-leto-03_integral}
\input{src/q_2022-podzim-03_posloupnosti}
\input{src/q_2023-podzim-04_geometricka-rada}
\input{src/q_2024-jaro-04_nespojitost-funkce}
\input{src/q_2024-podzim-03_integraly-primitivni-funkce}
\input{src/q_2025-podzim-03_zaklady-analyzy}
% KONEC-OTAZEK

\end{document}
